% Correspondance entre les relations injectées et les conclusions du chapitre d'analyse.
%
% CE FICHIER EST UNE ANNEXE, ET DEUX CONTRÔLES LE SUIVENT. `tests/test_provenance_des_chapitres.py`
% l'examine au même titre qu'un fichier de chapitre, et
% `tests/test_correspondance_relations_conclusions.py` y lit les lignes de tableau quand il lit les
% marques de conclusion dans le chapitre. Les deux portent leur liste de fichiers en clair.
%
% Le tableau occupait plus de deux pages du corps du rapport pour une matière que le lecteur
% consulte au cas par cas. Le chapitre garde à sa place ce qu'il établit, et renvoie ici.
%
% Ce fichier repose sur :
%   DOC: (aucun)
%   OBS: (aucun)
%   HYP: (aucun)
%   CHI: cible-delai-niveau-1, cible-delai-niveau-5, creances-part-au-dela-un-an, delai-rapport-extremes, doublons-part-mesuree, feries-programme-passages
%   CHI: flux-heure-pointe, flux-heures-vides-programme, hebdo-rapport-mesure, injecte-correlation-parametres, injecte-ferie-coefficient, injecte-hebdo-rapport
%   CHI: injecte-part-deux-derniers-niveaux, injecte-part-jour-meme, injecte-pente-absenteisme, injecte-ramadan-coefficient, injecte-taux-annulation, injecte-taux-doublons
%   CHI: part-jour-meme, qualite-colonne-vide, qualite-colonnes-incompletes, ramadan-baisse, recouvrement-taux-maximal, recouvrement-taux-minimal
%   CHI: relation-activites-ecart-positif, relation-inter-ensemble, taux-absenteisme-global, taux-annulation-global, urgences-part-consultation-ordinaire, urgences-part-retour-domicile
%   SER: (aucune)
%   SECTIONS: 0

\chapitreDAnnexe[Relations injectées et conclusions]{Correspondance entre relations injectées et conclusions}\label{annexe:correspondance}

Le jeu de données porte des relations injectées à dessein : un taux d'absentéisme fixé par
activité, une pente entre délai et absence, une répartition des niveaux de tri. Une conclusion du
chapitre~\ref{chap:activite} qui reproduit l'une d'elles ne découvre rien — elle montre que la
chaîne sait la retrouver.

Les deux tableaux qui suivent tiennent la correspondance dans les deux sens. Le premier va des
relations aux conclusions ; le second des conclusions à leur origine.

\textbf{Premier sens} : toute relation du registre des relations injectées est reprise par une
conclusion de ce chapitre, avec l'amplitude injectée en regard de l'amplitude mesurée, ou déclarée
non reprise avec son motif.

{\footnotesize
% Un tableau long, et non un tableau ordinaire : vingt et une lignes ne tiennent pas dans le
% reste d'une page, et un tableau ordinaire les reporte en bloc en laissant la page presque vide.
\begin{longtable}{@{}llp{0.255\linewidth}p{0.255\linewidth}@{}}
\toprule
Relation & Conclusion & Amplitude injectée & Amplitude mesurée \\
\midrule
\endfirsthead
\toprule
Relation & Conclusion & Amplitude injectée & Amplitude mesurée \\
\midrule
\endhead
\bottomrule
\endlastfoot
\relreprise{R-01}{C-11}{pente de \chiffre{injecte-pente-absenteisme} par jour, interne à la spécialité}{corrélation interne positive, écart de médiane positif sur \chiffre{relation-activites-ecart-positif} activités}
\relreprise{R-02}{C-09}{taux d'absentéisme fixé par activité}{\chiffre{taux-absenteisme-global}~\% en global, classement par activité retrouvé}
\relreprise{R-03}{C-07}{loi log-normale, médiane fixée par activité}{rapport de \chiffre{delai-rapport-extremes} entre les extrêmes, écart sous le jour entre sous-populations}
\relreprise{R-04}{C-02, C-13}{fenêtre de pointe posée, poids tous positifs aux urgences}{pointe à \chiffre{flux-heure-pointe} h, \chiffre{flux-heures-vides-programme} heures vides côté programmé, aucune aux urgences}
\relreprise{R-05}{C-03}{rapport de \chiffre{injecte-hebdo-rapport} entre le premier et le sixième jour}{\chiffre{hebdo-rapport-mesure}}
\relreprise{R-06}{C-04}{coefficient de \chiffre{injecte-ramadan-coefficient} pendant le Ramadan}{baisse de \chiffre{ramadan-baisse}~\%, un point sous le coefficient}
\relreprise{R-07}{C-05}{coefficient de \chiffre{injecte-ferie-coefficient} les jours fériés pour le programmé}{\chiffre{feries-programme-passages} passage programmé, urgences invariantes}
\relnonreprise{R-08}{aucune grandeur de ce chapitre ne ventile par régime de couverture : la dimension d'organisme payeur n'est référencée par aucun fait}
\relreprise{R-09}{C-14, C-16}{part de \chiffre{injecte-part-deux-derniers-niveaux} pour les deux niveaux aux cibles les plus longues}{\chiffre{urgences-part-consultation-ordinaire}~\%}
\relreprise{R-10}{C-14}{cibles de \chiffre{cible-delai-niveau-1} à \chiffre{cible-delai-niveau-5} min, respect fixé}{chaque médiane sous sa cible, chaque neuvième décile au-dessus}
\relreprise{R-11}{C-15}{répartition fixée des cinq issues}{\chiffre{urgences-part-retour-domicile}~\% de retour à domicile, transfert devant hospitalisation}
\relnonreprise{R-12}{aucune conclusion de ce chapitre ne croise le champ d'origine avec le niveau de tri ; la relation est traitée au chapitre de conception du jeu de données}
\relnonreprise{R-13}{aucune colonne d'âge ni de tranche d'âge n'est calculée dans les couches aval : le chapitre n'a aucune grandeur démographique à analyser}
\relreprise{R-14}{C-22}{taux de doublons de \chiffre{injecte-taux-doublons}}{\chiffre{doublons-part-mesuree}~\%}
\relreprise{R-15}{C-20}{taux de champs manquants fixé par colonne}{\chiffre{qualite-colonnes-incompletes} colonnes incomplètes, dont \chiffre{qualite-colonne-vide} entièrement vide}
\relreprise{R-16}{C-09}{taux d'annulation de \chiffre{injecte-taux-annulation}}{\chiffre{taux-annulation-global}~\%}
\relreprise{R-17}{C-08}{part de \chiffre{injecte-part-jour-meme} pour le jour même}{\chiffre{part-jour-meme}~\%, l'excès venant de la marge nulle}
\relnonreprise{R-18}{aucune page ne donne la part des patients adressés à voir, et ce chapitre n'ouvre aucune section sur l'adressage}
\relnonreprise{R-19}{le laboratoire n'entre dans aucune des six sections de ce chapitre ; le ratio est traité au chapitre de conception du jeu de données}
\relreprise{R-20}{C-18, C-19}{taux de facturation et de recouvrement fixés par type}{de \chiffre{recouvrement-taux-minimal} à \chiffre{recouvrement-taux-maximal}~\% selon le débiteur ; \chiffre{creances-part-au-dela-un-an}~\% du restant dû a un an}
\relreprise{R-21}{C-11}{corrélation de \chiffre{injecte-correlation-parametres} entre les deux tables de paramètres}{\chiffre{relation-inter-ensemble}}
\end{longtable}
}

\textbf{Second sens} : toute conclusion renvoie à une relation injectée, ou est déclarée comme n'en
venant pas — auquel cas elle est un effet des autres choix du générateur et \textbf{n'établit
rien}.

\begin{center}
\footnotesize
\begin{tabular}{@{}lp{0.72\linewidth}@{}}
\toprule
Conclusion & Ce dont elle vient \\
\midrule
\conclusionhors{C-01}{d'aucune relation : effet des paramètres de volumétrie et du calendrier}
\conclusionhors{C-06}{d'aucune relation : effet du mécanisme de plafonnement du délai}
\conclusionhors{C-10}{d'aucune relation : R-16 déclare que la distinction relève du modèle de données. C'est ce que la chaîne établit vraiment}
\conclusionhors{C-12}{d'aucune relation : propriété du mécanisme de plafonnement, mesurée sur le complément}
\conclusionhors{C-17}{d'aucune relation : effet conjoint des tarifs et de la clé de couverture}
\conclusionhors{C-21}{d'aucune relation : effet des défauts de forme semés à l'ingestion}
\bottomrule
\end{tabular}
\end{center}
